\section{Introduction}

In this chapter, we want to see how the basic objects and constructions provided by Euclidean geometry can be used to build a Cartesian coordinate system.

Points, lines, segments, circles, perpendicularity, and parallelism form a geometric world in which we can construct objects and study their relations. Choosing axes, an origin, directions, and a unit allows us to describe that same world using numbers.

This transition changes the way geometry can be practised. Geometric objects acquire algebraic descriptions, while position, distance, and mutual relations can be studied through calculation. A Cartesian coordinate system does not replace geometry; it gives geometry a new language and opens access to a broad range of algebraic methods.
